% Remaster, 6 October 2026, after the supervisor's comments on the abstract
% (Overleaf, 2 October).  Journal limit 250 words.  Structure: context, gap,
% what we did, three results, implication.
Agentic systems built on large language models are increasingly used to
automate machine learning. Where their accuracy comes
from is unclear: the language model, the harness around it, or the expert
instruction file (an agent skill) it is given. We separate the three. All
experiments use one open-weight 31B language model and three harnesses: FEDOT.LLM (a fixed-stage workflow over an AutoML library), AutoDS-Tools (a
six-agent system) and Terminus-2 (a minimal shell agent). Accuracy is measured against the eighteen
hyperparameter-tuned methods published with the TabReD benchmark, on a
per-task scale from 0 (weakest published method) to 1 (strongest). First,
without the AutoDS-Tools instruction file the mean scores of the three
harnesses are close ($\nTerminus$ to $\nFedot$) and no pairwise difference
is significant. Second, accuracy follows who builds and tunes the model. An
agent writing its own code fits one untuned gradient-boosting model. The
instruction file raises AutoDS-Tools from $\nAutoDSOff$ to $\nAutoDS$, the
level of tuned LightGBM, on all eight tasks, through the part that names an
AutoML library. The strongest published method ($\nBest$) remains out of
reach. Third, the file and the language model change reliability, and the
direction depends on the harness. In another harness the same file ends
\ktBothLost{} of 24 trials without a submission, and a larger model removes
such trials in one harness and adds them in the other. In practice, the
harness decides reliability and run time, an AutoML library decides
accuracy, and an instruction file must be tested with its harness.
